\documentclass{article}
\usepackage[utf8]{inputenc}
\usepackage{a4wide}

\begin{document}

\section*{เดินเที่ยวนิวยอร์ก}

ถนนในแมนแฮตตันเมืองนิวยอร์กมีผังเมืองเป็นตารางเกือบสมบูรณ์

กำหนดอาร์เรย์ของจุดชุดหนึ่ง\texttt{xi,\ yi}เป็นพิกัดสมมติแทนสถานที่ในแมนแฮตตันที่

ดร.ต๋อยจะต้องเดินทางไปตามเส้นทาง\textbf{( ไปทางเหนือ, ใต้, ตะวันออก, หรือตะวันตก

)}ระยะทางจะ\textbf{คิดแบบระยะทางแมนแฮตตัน ( Manhattan distance

)}สำหรับเชื่อมต่อจุด\texttt{xi,\ yi}กับ\texttt{xj,\ yj}ซึ่งคำนวณได้จาก{}{\texttt{\textbar{}}}\texttt{xi–xj}{\texttt{\textbar{}\ +\ \textbar{}}}\texttt{yi–yj}{\texttt{\textbar{}}}{}โดยที่\texttt{\textbar{}v\textbar{}}หมายถึงค่าสัมบูรณ์ของ\texttt{v}



\textbf{\ul{ให้หาระยะทางสั้นที่สุดที่

ดร.ต๋อยจะหาเส้นทางที่เชื่อมต่อทุกจุด}\ul{ที่กำหนดให้}}นั่นคือจะต้องมีเส้นเชื่อม 1

เส้นระหว่างจุดสองจุด



\hfill\break



\textbf{ตัวอย่าง}



\begin{quote}

\textbf{Input:}



5



0 0



2 2



3 10



5 2



7 0



\textbf{Output:}



20

\end{quote}



\hfill\break



\includegraphics[width=1.875in,height=2.39583in,alt={A graph with black dots Description automatically generated}]{https://lh7-us.googleusercontent.com/l2SHVHbo3qXHzWJNrHH8_ZNTil1A_8DfAW4r9nvP58EAwXblA4cMld7KYNKCq1XjDPUuzfpvNu5r_lxHo7rappv7bLKjGgJyZ-fAB6vOqcl7MkwpIo8F4KTfbMiQKaB6Bhr7yYXbBIcPF-zyYbOmfA}\includegraphics[width=1.875in,height=2.39583in,alt={A graph of a cost Description automatically generated}]{https://lh7-us.googleusercontent.com/DVKIPkGBewkgP6Tv6BeqybSoKxjxNxRUfNq-3pVu9rJyLY8MZ3DYcfkJTE1ETIVxZYxOSxRFoGSdVvRIzFD4pE2SYiVeqLYHIfFpf7TmbtiVGJkv5ZCur6EATVj-jU0vZuAfpsiFKs7_3b7xe1jhjA}\\



\subsection*{Input}
เป็นบรรทัดแรก จำนวนจุด\texttt{N}



บรรทัดที่ 2 จนถึง บรรทัดที่\texttt{N\ +\ 1}เป็นพิกัด\texttt{xi\ yi}คั่นกลางด้วยช่องว่าง

1 ช่อง{โดย}{พิกัดทุกค่าเป็นจำนวนเต็มบวก}



\subsection*{Output}
เป็นจำนวนเต็มบวกแทนระยะทางแมนแฮตตันที่\ul{\textbf{ต่ำสุด}}ที่เชื่อมต่อทุกจุดที่กำหนด\\



\end{document}
